\subsection{深度与正规性}

设 A 为 Noether 环，p 为 A 的素理想。有 \(\text{prof}(A_{\mathfrak{p}}) \leqslant \text{ht}(\mathfrak{p})\) （第 4 节, 定理 2 的推论 2）。此外若 A 约化，则局部环 \(A_{p}\) 约化 (II, § 2, 第 6 节, 命题 17)，由此蕴含：

\begin{enumerate}
\def\labelenumi{\alph{enumi})}
\item
  若 ht(p) = 0，则 A \(_{p}\) 是域；
\item
  若 \(\operatorname{ht}(\mathfrak{p}) \geqslant 1\) ，则 \(\operatorname{prof}(A_{\mathfrak{p}}) \geqslant 1\) （ \(n^{\circ}\) 1, 注 3）。
\end{enumerate}

反之：

命题 15.--- 设 A 为满足下列两个条件的 Noether 环：

\begin{enumerate}
\def\labelenumi{(\roman{enumi})}
\item
  对 A 的每个极小素理想 \(\mathfrak{p}\) ，环 \(A_{\mathfrak{p}}\) 约化；
\item
  对 A 的每个素理想 \(\mathfrak{p}\) （高度 \(\geqslant 1\) 者），有 \(\operatorname{prof}(A_{\mathfrak{p}}) \geqslant 1\) 。则 A 约化。
\end{enumerate}

设 \(\mathfrak{n}\) 为 A 的诣零根。对 A 的每个极小素理想 \(\mathfrak{p}\) ，由 (i) 有 \(\mathfrak{n}_{\mathfrak{p}} = 0\) ，也就是 \(\mathfrak{p} \notin \operatorname{Supp}(\mathfrak{n})\) ，更有 \(\mathfrak{p} \notin \operatorname{Ass}_{\Lambda}(\mathfrak{n})\) 。对每个 \(\mathfrak{p} \in \operatorname{Spec}(\mathbf{A})\) （高度 \(\geqslant 1\) 者），由 有 \(\mathfrak{pA}_{\mathfrak{p}} \notin \operatorname{Ass}_{\mathrm{A}_{\mathfrak{p}}}(\mathrm{A}_{\mathfrak{p}})\) ，更有 \(\mathfrak{pA}_{\mathfrak{p}} \notin \operatorname{Ass}_{\mathrm{A}_{\mathfrak{p}}}(\mathfrak{n}_{\mathfrak{p}})\) ，从而 \(\mathfrak{p} \notin \operatorname{Ass}_{\mathrm{A}}(\mathfrak{n})\) (IV, § 1, 第 2 节, 命题 5)。于是集合 \(\operatorname{Ass}_{\mathrm{A}}(\mathfrak{n})\) 为空，这蕴含 \(\mathfrak{n}\) 为零。

命题 16.--- 设 A 为整闭的 Noether 环，J 为 A 的高度 ≥ 2 的理想，M 为自反的有限生成 A-模。则 prof \(_{A}\) (J; M) ≥ 2。

选取 A 的分式域上的一个有限维向量空间 V，以及 V 中同构于 M 的一个格 N (VII, § 4, 第 2 节, 注 1)。V/N 的相伴素理想高度为 1（同上引文, 定理 2），故不包含 J；由第 1 节的注 2，这蕴含 \(\operatorname{prof}_{\mathrm{A}}(\mathrm{J};\mathrm{V}/\mathrm{N}) \geqslant 1\) 。另一方面，A-模 V 是可除的且无挠，故为内射的 (A, X, 第 17 页, 命题 10 的推论 2)，这蕴含 \(\operatorname{prof}_{\mathrm{A}}(\mathrm{J};\mathrm{V}) = +\infty\) 。于是不等式 \(\operatorname{prof}_{\mathrm{A}}(\mathrm{J};\mathrm{N}) \geqslant 2\) 由第 1 节的命题 1 得出。

推论.--- 整闭且维数 ≥ 2 的 Noether 局部环深度 ≥ 2。

设 A 为正规环（第 8 节），p 为 A 的素理想。则：

\begin{enumerate}
\def\labelenumi{\alph{enumi})}
\item
  若 ht(p) = 0，则 A \(_{p}\) 是域；
\item
  若 \(\operatorname{ht}(\mathfrak{p}) = 1\) ，则 \(A_{p}\) 是离散赋值环 (VII, § 1, 第 7 节, 命题 11)；
\item
  若 \(\operatorname{ht}(\mathfrak{p}) \geqslant 2\)，则 \(\operatorname{prof}(A_{\mathfrak{p}}) \geqslant 2\) （命题 16 的推论）。
\end{enumerate}

反之：

定理 4 (Serre).--- 设 A 为满足下列条件的 Noether 环：

\begin{enumerate}
\def\labelenumi{(\roman{enumi})}
\item
  对 A 的每个极小素理想 p，环 A \(_{p}\) 约化；
\item
  对 A 的每个高度 1 的素理想 p，环 A \(_{p}\) 整闭；
\item
  对 A 的每个素理想 \(\mathfrak{p}\) （高度 \(\geqslant 2\) 者），有 \(\operatorname{prof}(A_{\mathfrak{p}}) \geqslant 2\) 。
\end{enumerate}

则 A 正规。

目标是证明环 \(A_{\mathfrak{p}}\) 对每个 \(\mathfrak{p} \in \operatorname{Spec}(A)\) 整闭，我们对 \(\operatorname{ht}(\mathfrak{p})\) 作归纳来证明这一点。对 \(\operatorname{ht}(\mathfrak{p}) \leqslant 1\)，这由假设 与 得出。于是设 \(\operatorname{ht}(\mathfrak{p}) \geqslant 2\) ，且 \(A_{\mathfrak{q}}\) 对 A 的每个素理想 \(\mathfrak{q}\) （满足 \(\operatorname{ht}(\mathfrak{q}) < \operatorname{ht}(\mathfrak{p})\) 者）整闭。由假设，有 \(\operatorname{prof}(A_{\mathfrak{p}}) \geqslant 2\) 。由归纳假设及第 9 节的定理 3 的推论 3——应用于环 \(A_{\mathfrak{p}}\) 及闭子集 \(\{\mathfrak{p}A_{\mathfrak{p}}\}\) （ \(\operatorname{Spec}(A_{\mathfrak{p}})\) 的闭子集）——环 \(A_{\mathfrak{p}}\) 是整环。设 K 为它的分式域，设 B 为 K 的一个包含 \(A_{\mathfrak{p}}\) 且在 \(A_{\mathfrak{p}}\) 上有限的子环。剩下要证明 B 等于 \(A_{\mathfrak{p}}\) 。用 i 表示从 \(A_{\mathfrak{p}}\) 到 B 的典范单射。由于 B 包含于 K，它是无挠的 \(A_{\mathfrak{p}}\) -模，故深度 \(\geqslant 1\) 。此外，对每个素理想 \(\mathfrak{q}\) （ \(A_{\mathfrak{p}}\) 的素理想，异于 \(\mathfrak{p}A_{\mathfrak{p}}\) 者），同态 \(i_{\mathfrak{q}}: (A_{\mathfrak{p}})_{\mathfrak{q}} \to B_{\mathfrak{q}}\) 是双射，因为 \(A_{\mathfrak{q}}\) 整闭。由第 9 节的引理 4——应用于闭子集 \(F = \{\mathfrak{p}A_{\mathfrak{p}}\}\) （ \(\operatorname{Spec}(A_{\mathfrak{p}})\) 的闭子集）——同态 i 是双射，这就完成了定理的证明。

注 1.--- 定理 4 的一个方便的形式如下：设 A 为这样的 Noether 环：对 A 的每个素理想 p，环 A \(_{p}\) 整闭或深度 ≥ 2；则 A 正规。事实上，我们来验证定理 4 的假设。设 p ∈ Spec(A)。若 ht(p) ≤ 1，则有 prof(A \(_{p}\) ) ≤ 1，故 A \(_{p}\) 整闭。若 ht(p) ≥ 2，则环 A \(_{p}\) 或者深度 ≥ 2，或者整闭，后者同样蕴含深度(A \(_{p}\) ) ≥ 2（第 1 节, 命题 1 的推论 2），由此得所要的结论。类似地可验证下述命题：设 A 为这样的 Noether 环：对 A 的每个素理想 p，环 A \(_{p}\) 约化或深度 ≥ 1；则 A 约化。

推论 1.--- 设 \(A\) 为 Noether 环，F 为 \(\operatorname{Spec}(A)\) 的闭子集，U 为其补开集。假设 \(\operatorname{prof}_{\mathrm{F}}(\mathrm{A})\) 为 \(\geqslant 2\) （相应地为 \(\geqslant 1\) ），且对每个 \(\mathfrak{p} \in \mathrm{U}\) ，环 \(\mathrm{A}_{\mathfrak{p}}\) 整闭（相应地约化）。则 A 正规（相应地约化）。

对每个 \(\mathfrak{p} \in \mathrm{F}\) ，有 \(\operatorname{prof}(\mathrm{A}_{\mathfrak{p}}) \geqslant \operatorname{prof}_{\mathrm{F}}(\mathrm{A})\) （ \(n^{\circ}\) 5, 命题 8）；因此只需应用前一个注。

推论 2.-- 设 ρ : A → B 为使 B 成为平坦 A-模的 Noether 环之间的同态。

\begin{enumerate}
\def\labelenumi{\alph{enumi})}
\item
  若 B 正规且在 A 上忠实平坦，则 A 正规。
\item
  设 A 正规，且环 \(\kappa(\mathfrak{p})\otimes_{\mathrm{A}}\mathrm{B}\) 当 \(\mathfrak{p}\) 为 A 的极小素理想时正规、当 \(\mathfrak{p}\) 为高度 1 的素理想时约化。则环 B 正规。
\item
  设 B 正规且在 A 上忠实平坦。则 \(\rho\) 是单射 (I, § 3, 第 5 节, 命题 9) 且 B 约化，故 A 约化。设 S 为 A 中非零因子构成的集合。由于 B 在 A 上平坦，\(\rho(S)\) 由 B 中的非零因子构成，故 B 在 \(\rho(S)^{-1}B\) 中整闭。设 \(x\) 为 \(S^{-1}A\) 中在 A 上整的元素；则元素 \(x \otimes 1_B\) ——环 \(S^{-1}A \otimes_A B\) （它等同于 \(\rho(S)^{-1}B\) ）的元素——在 B 上整，故属于 B。由于 B 在 A 上忠实平坦，有 \(x \in A\) (I, § 3, 第 5 节, 命题 10, (ii))，故 A 正规。
\item
  在 b) 的假设下，由注 1 只需证明：对 B 的每个素理想 \(\mathfrak{q}\) ，局部环 \(B_{\mathfrak{q}}\) 正规或深度 \(\geqslant 2\) 。用 \(\mathfrak{p}\) 表示 A 的素理想 \(\rho^{-1}(\mathfrak{q})\) ；局部同态 \(A_{\mathfrak{p}} \to B_{\mathfrak{q}}\) ——由 \(\rho\) 诱导——使 \(B_{\mathfrak{q}}\) 成为忠实平坦的 \(A_{\mathfrak{p}}\) -模 (I, § 3, 第 5 节, 命题 9)。我们区分四种情形：
\end{enumerate}

\begin{enumerate}
\def\labelenumi{\arabic{enumi})}
\item
  若 \(\operatorname{ht}(\mathfrak{p}) = 0\) ，则 \(A_{\mathfrak{p}}\) 是域，等于 \(\kappa(\mathfrak{p})\) ；环 \(A_{\mathfrak{p}} \otimes_{A} B\) 按假设正规。\(B_{\mathfrak{q}}\) 亦然，它是其一个分式环。
\item
  若 \(\operatorname{ht}(\mathfrak{p}) = 1\) 且 \(\operatorname{ht}(\mathfrak{q}) \leqslant 1\) ，则 \(A_{\mathfrak{p}}\) 是离散赋值环；设 \(\pi\) 为 \(A_{\mathfrak{p}}\) 的一个单值化元。由于 \(B_{\mathfrak{q}}\) 在 \(A_{\mathfrak{p}}\) 上忠实平坦，元素 \(\pi 1_{B_{\mathfrak{q}}}\) ——\(B_{\mathfrak{q}}\) 的元素——不是零因子，故局部环 \(B_{\mathfrak{q}}/\pi B_{\mathfrak{q}}\) 的维数为 0 (VIII, § 3, 第 1 节, 命题 1 的推论 2)。它是约化的，因为它是约化环 \(\kappa(\mathfrak{p}) \otimes_{A} B\) 的一个分式环，从而它是域。因此 \(B_{\mathfrak{q}}\) 是离散赋值环 (VI, § 1, 第 4 节, 命题 2)，故整闭。
\item
  \(\operatorname{ht}(\mathfrak{p}) = 1\) 且 \(\operatorname{ht}(\mathfrak{q}) \geqslant 2\) 。则由关系式
\end{enumerate}

\[
\dim (B _ {\mathfrak {q}}) = \dim (A _ {\mathfrak {p}}) + \dim (\kappa (\mathfrak {p}) \otimes_ {A} B _ {\mathfrak {q}})
\]

(VIII, § 3, 第 4 节, 命题 7 的推论 1)，环 \(\kappa(\mathfrak{p})\otimes_{\mathrm{A}}\mathrm{B}_{\mathfrak{q}}\) 的维数 \(\geqslant 1\) 。它按假设约化，故深度 \(\geqslant 1\) （第 1 节, 注 3）。于是有（第 6 节, 命题 11 的推论）

\[
\operatorname{prof} \left(\mathrm{B} _ {\mathfrak {q}}\right) = \operatorname{prof} \left(\Lambda_ {\mathfrak {p}}\right) + \operatorname{prof} \left(\kappa (\mathfrak {p}) \otimes_ {\mathrm{A}} \mathrm{B} _ {\mathfrak {q}}\right) \geqslant 1 + 1 = 2.
\]

\begin{enumerate}
\def\labelenumi{\arabic{enumi})}
\setcounter{enumi}{3}
\tightlist
\item
  \(\operatorname{ht}(\mathfrak{p}) \geqslant 2\) 。由于 \(A_{p}\) 的深度 \(\geqslant 2\) （命题 16 的推论），\(B_{q}\) 由同上引文亦然。
\end{enumerate}

推论 3.--- 设 \(\rho: A \to B\) 为 Noether 环之间的同态。设 \(B\) 是平坦 A-模，\(A\) 正规，且 \(\kappa(\mathfrak{p}) \otimes_A B\) 对所有 \(\mathfrak{p} \in \operatorname{Spec}(A)\) 正规。则 \(B\) 正规。

